\subsection{Homogeneous Systems}
\hrulefill

\paragraph*{Definition} A system is called \textbf{homogeneous} if all the constant terms are zero. In other words, a system of equations is homogeneous if it can be written in the form
\[\begin{cases}
    x-5y = 0 \\
    2x+3y+z = 0 \\
\end{cases}
\]
The soluttion set of a homogeneous system always contains the \textbf{trivial solution}, which is the solution where all variables are equal to zero. 
For example, in the system above, the trivial solution is \(x=0, y=0, z=0\).  Any other solution is called a \textbf{nontrivial solution}.

\paragraph*{Example} Find any non-trivial solution
\begin{align*}
    &\begin{bmatrix}
        1 & -1 & 2 & -1 & | & 0 \\
        2 & 2 & 0 & 1 & | & 0 \\
        3 & 1 & 2 & -1 & | & 0
    \end{bmatrix} \xmapsto{R_2 \mapsto R_2 - 2R_1}\begin{bmatrix}
        1 & -1 & 2 & -1 & | & 0 \\
        0 & 4 & -4 & 3 & | & 0 \\
        3 & 1 & 2 & -1 & | & 0
    \end{bmatrix} \\
    \xmapsto{R_3 \mapsto R_3 - 3R_1} &\begin{bmatrix}
        1 & -1 & 2 & -1 & | & 0 \\
        0 & 4 & -4 & 3 & | & 0 \\
        0 & 4 & -4 & 2 & | & 0
    \end{bmatrix} \xmapsto{R_3 \mapsto R_3 - R_2} \begin{bmatrix}
        1 & -1 & 2 & -1 & | & 0 \\
        0 & 4 & -4 & 3 & | & 0 \\
        0 & 0 & 0 & -1 & | & 0
    \end{bmatrix} \\
    \xmapsto{R_1 \mapsto R_1 - R_3} &\begin{bmatrix} 
        1 & -1 & 2 & 0 & | & 0 \\
        0 & 4 & -4 & 3 & | & 0 \\
        0 & 0 & 0 & -1 & | & 0
    \end{bmatrix} \xmapsto{R_2 \mapsto R_2 + 3R_3} \begin{bmatrix}
        1 & -1 & 2 & 0 & | & 0 \\
        0 & 4 & -4 & 0 & | & 0 \\
        0 & 0 & 0 & -1 & | & 0
    \end{bmatrix} \\
    \xmapsto{R_3 \mapsto -R_3} &\begin{bmatrix}
        1 & -1 & 2 & 0 & | & 0 \\
        0 & 4 & -4 & 0 & | & 0 \\
        0 & 0 & 0 & 1 & | & 0
    \end{bmatrix} \xmapsto{R_2 \mapsto \frac{1}{4}R_2} \begin{bmatrix}
        1 & -1 & 2 & 0 & | & 0 \\
        0 & 1 & -1 & 0 & | & 0 \\
        0 & 0 & 0 & 1 & | & 0
    \end{bmatrix} \\
    \xmapsto{R_1 \mapsto R_1 + R_2} &\begin{bmatrix}
        1 & 0 & 1 & 0 & | & 0 \\
        0 & 1 & -1 & 0 & | & 0 \\
        0 & 0 & 0 & 1 & | & 0
    \end{bmatrix} \implies \begin{cases}
        x_1 + x_3 = 0 \\
        x_2 - x_3 = 0 \\
        x_4 = 0
    \end{cases} = \begin{cases}
        x_1 = -x_3 \\
        x_2 = x_3 \\
        x_4 = 0
    \end{cases} \\
    x_3 = S \implies &\begin{cases}
        x_1 = -S \\
        x_2 = S \\
        x_3 = S \\
        x_4 = 0
    \end{cases} S \in \mathbb{R} \qquad S = 1 \implies \begin{cases}
        x_1 = -1 \\
        x_2 = 1 \\
        x_3 = 1 \\
        x_4 = 0
    \end{cases}
\end{align*}

\paragraph*{Theorem} If a homogeneous system of linear equations has ore variables than equations, 
then it has infinitly many non-trivial solutions. For some matrix $A$ with $m$ rows and $n$ parameters
 $\text{rank}(A) \leq m < n$ implies that the solutions involved $n-rank(A) \geq 1$ parameters, so $x_i = S$ and 
 we can choose any value for $S$ other than zero to get a non-trivial solution.
\paragraph*{Note} If a homogeneous system has a nontrivial solution, it need not have more variables than equations.
\paragraph*{Note} If a homogeneous system has a non-trivial solution, then it has infinitely many non-trivial solutions. If $x = S$ is a non-trivial solution, then so is $x = kS$ for any scalar $k \neq 0$.

\hrulefill

\subsubsection*{Vectors}
\paragraph*{Definition} A \textbf{vector} in $\mathbb{R}^n$ is an ordered $n$-tuple of real numbers. We can represent a vector as a column matrix:
\[\vec{x} = \begin{bmatrix}
    x_1 \\
    x_2 \\
    \vdots \\
    x_n
\end{bmatrix} = (x_1, x_2, \ldots, x_n) = \langle x_1, x_2, \ldots, x_n \rangle\]
where $x_1, x_2, \ldots, x_n$ are the components of the vector $\vec{x}$. In this class, vectors should always be 
represented as column matrices unless otherwise specified.

\paragraph*{Definition} A vector of length $n$ is an ordered n-tuple, and the set of all vectors of length $n$ is denoted by $\mathbb{R}^n$.
In other words, $\vec{x} \in \mathbb{R}^n$ if $\vec{x}$ has $n$ components.

\paragraph*{Vector Addition} Vector addition can only be performed on vectors of the same length. If $\vec{x}, \vec{y} \in \mathbb{R}^n$, then their sum is given by
\[\vec{x} + \vec{y} = \begin{bmatrix}
    x_1 \\
    x_2 \\
    \vdots \\
    x_n
\end{bmatrix} + \begin{bmatrix}
    y_1 \\
    y_2 \\
    \vdots \\
    y_n
\end{bmatrix} = \begin{bmatrix}
    x_1 + y_1 \\
    x_2 + y_2 \\
    \vdots \\
    x_n + y_n
\end{bmatrix}\]

\paragraph*{Scalar Multiplication} If $\vec{x} \in \mathbb{R}^n$ and $c$ is a scalar, then the scalar multiplication of $c$ and $\vec{x}$ is given by
\[c\vec{x} = c \begin{bmatrix}
    x_1 \\
    x_2 \\
    \vdots \\
    x_n
\end{bmatrix} = \begin{bmatrix}
    cx_1 \\
    cx_2 \\
    \vdots \\
    cx_n
\end{bmatrix}\]

\paragraph*{Linear Combination}
A sum of scalar multiples of several vectors or columns is called a \textbf{linear combination} of the vectors or columns.
For example, if $\vec{u}, \vec{v}, \vec{w} \in \mathbb{R}^n$ and $a, b, c$ are scalars, then
\[a\vec{u} + b\vec{v} + c\vec{w}\]
is a linear combination of the vectors $\vec{u}, \vec{v}, \vec{w}$.
\pagebreak

\paragraph*{Example} Let $\vec{x} = \begin{bmatrix} 0 \\ 1 \\ 0 \end{bmatrix}, \vec{y} = \begin{bmatrix} 2 \\ 1 \\ 0 \end{bmatrix},
\vec{z} = \begin{bmatrix} 3 \\ 1 \\ 1 \end{bmatrix}, \vec{v} = \begin{bmatrix} 0 \\ -1 \\ 2 \end{bmatrix}, \vec{w} = \begin{bmatrix}
1 \\ 1 \\ 1 \end{bmatrix}$. \\ a. Find $2\vec{x} -\vec{y} + 3\vec{z}$, b. Determine if $\vec{v}$ is a linear combination
of $\vec{x}, \vec{y}, \vec{z}$, and c. Determine if $\vec{w}$ is a linear combination of $\vec{x}, \vec{y}, \vec{z}$.

\begin{align*}
    a. \quad & 2\vec{x} - \vec{y} + 3\vec{z} = 2\begin{bmatrix} 0 \\ 1 \\ 0 \end{bmatrix} - \begin{bmatrix} 2 \\ 1 \\ 0 \end{bmatrix} + 3\begin{bmatrix} 3 \\ 1 \\ 1 \end{bmatrix} = \begin{bmatrix}
        9 \\
        5 \\
        3
    \end{bmatrix} \\
    b. \quad & \exists r, s, t \in \mathbb{R} \ni r\vec{x} + s\vec{y} + t\vec{z} = \vec{v} \implies \begin{bmatrix} r \\ 0 \\ r \end{bmatrix} + \begin{bmatrix} 2s \\ s \\ 0 \end{bmatrix} + \begin{bmatrix} 3t \\ t \\ t \end{bmatrix} = \begin{bmatrix} 0 \\ -1 \\ 2 \end{bmatrix} \\
    \implies & \begin{cases}
        r + 2s + 3t = 0 \\
        s + t = -1 \\
        r + t = 2
    \end{cases} \implies \begin{bmatrix}
        1 & 2 & 3 & | & 0 \\
        0 & 1 & 1 & | & -1 \\
        1 & 0 & 1 & | & 2
    \end{bmatrix} \xmapsto{RREF} \begin{bmatrix}
        1 & 0 & 1 & | & 2 \\
        0 & 1 & 1 & | & -1 \\
        0 & 0 & 0 & | & 0
    \end{bmatrix} \\
    \implies & \begin{cases}
        r + t = 2 \\
        s + t = -1
    \end{cases} \implies \begin{cases}
        r = 2 - q \\
        s = -1 - q \\
        t = q
    \end{cases} t \in \mathbb{R} \quad q = 0 \implies \begin{cases}
        r = 2 \\
        s = -1 \\
        t = 0
    \end{cases} \implies \vec{v} = 2\vec{x} - \vec{y} \\
    c. \quad & \exists r, s, t \in \mathbb{R} \ni r\vec{x} + s\vec{y} + t\vec{z} = \vec{w} \implies \begin{bmatrix}
        1, 2, 3 & | & 1 \\
        0, 1, 1 & | & 1 \\
        1, 0, 1 & | & 1
    \end{bmatrix} \xmapsto{RREF} \begin{bmatrix}
        1 & 0 & 1 & | & 1 \\
        0 & 1 & 1 & | & 1 \\
        0 & 0 & 0 & | & -1
    \end{bmatrix} \\
    \implies & \text{No solution, so $\vec{w}$ is not a linear combination of $\vec{x}, \vec{y}, \vec{z}$}
\end{align*} 

\paragraph*{Note} To find if $\vec{x}$ is a linear combinition of $\vec{a_1}, \vec{a_2}, \ldots, \vec{a_n}$, we can form the augmented matrix
\[[\vec{a_1} \ \vec{a_2} \ \ldots \ \vec{a_n} \ | \ \vec{x}]\]
and row reduce to see if there is a solution. Linear combinations are a good way to write the general solution to a homogeneous system.

\paragraph*{Theorem} Any linear combination of solutions to a homogeneous system is also a solution to the system.
\paragraph*{Example} Let $\vec{p} = \begin{bmatrix} p_1 \\ p_2 \end{bmatrix}$ and $\vec{q} = \begin{bmatrix} q_1 \\ q_2 \end{bmatrix}$ be solutions to the homogeneous system $a_1x_1 + a_2x_2 = 0$, hence:
\begin{align*}
    \implies &\exists s, t \in \mathbb{R} \ni \begin{cases}
        a_1p_1 + a_2p_2 = 0 \\
        a_1q_1 + a_2q_2 = 0
    \end{cases} = s\vec{p} + t\vec{q} = \begin{bmatrix}
        sp_1 + tq_1 \\
        sp_2 + tq_2
    \end{bmatrix} \text{ is a solution} \\
    &a_1(sp_1 + tq_1) + a_2(sp_2 + tq_2) = 0 \implies s \cdot 0 + t \cdot 0 = 0
\end{align*}

\pagebreak

\paragraph*{Example} Solve the homogeneous system w/ coefficient matrix $A = \begin{bmatrix}
    1 & -2 & 3 & -2 \\
    -3 & 6 & 1 & 0 \\
    -2 & 4 & 4 & -2
\end{bmatrix}$
\begin{align*}
    RREF(A) = &\begin{bmatrix}
        1 & -2 & 0 & -\frac{1}{5} & | & 0 \\
        0 & 0 & 1 & -\frac{3}{5} & | & 0 \\
        0 & 0 & 0 & 0 & | & 0
    \end{bmatrix} \quad \text{Rank}(A) = 2 \implies \text{4 - 2 = 2 parameters} \\
    \implies &\begin{cases}
        x_1 = 2s + \frac{1}{5}t \\
        x_2 = s \\
        x_3 = \frac{3}{5}t \\
        x_4 = t
    \end{cases} = \begin{bmatrix}
        2s + \frac{1}{5}t \\
        s \\
        \frac{3}{5}t \\
        t
    \end{bmatrix} = s\begin{bmatrix}
        2 \\
        1 \\
        0 \\
        0
    \end{bmatrix} + t\begin{bmatrix}
        \frac{1}{5} \\
        0 \\
        \frac{3}{5} \\
        1
    \end{bmatrix} \text{ is the general solution},\\
     \text{so our basic solutions are} \begin{bmatrix}
        2 \\
        1 \\
        0 \\
        0
    \end{bmatrix} \text{ and } \begin{bmatrix}
        \frac{1}{5} \\
        0 \\
        \frac{3}{5} \\
        1
    \end{bmatrix}
\end{align*}

\paragraph*{Note} Every solution to a homogeneous system can be written as a linear combination of the basic solutions.